\section*{Глава 1. Анализ предметной области}
\addcontentsline{toc}{section}{Глава 1. Анализ предметной области}
\stepcounter{section}

\subsection{Обзор существующих платформ для заказа еды и бронирования столиков}

На современном рынке цифровых сервисов для ресторанного бизнеса можно выделить несколько категорий решений: платформы доставки еды, системы онлайн-бронирования и комплексные решения для автоматизации кафе и ресторанов. Большинство существующих платформ ориентированы либо на доставку еды, либо на внутреннее управление заведением, редко объединяя весь необходимый функционал в единой системе.

Для определения функциональных особенностей разрабатываемого приложения был проведен сравнительный анализ популярных сервисов, представленный в таблице~\ref{tab:competitors}.

\begin{table}[ht!]
\centering
\caption{Сравнительный анализ существующих платформ}
\label{tab:competitors}
\begin{tabular}{|
>{\raggedright\arraybackslash}p{3.5cm}|
>{\centering\arraybackslash}p{2.4cm}|
>{\centering\arraybackslash}p{3.1cm}|
>{\centering\arraybackslash}p{3.1cm}|
>{\centering\arraybackslash}p{2.8cm}|}
\hline
\textbf{Платформа} &
\textbf{Онлайн-заказ еды} &
\textbf{Бронирование столиков} &
\textbf{Предзаказ блюд к брони} &
\textbf{Управление персоналом} \\ \hline

	Яндекс Еда &
	Да &
	Нет &
	Нет &
	Нет \\ \hline
	
	Delivery Club &
	Да &
	Нет &
	Нет &
	Нет \\ \hline
	
	RestoCRM &
	Частично &
	Частично &
	Нет &
	Да \\ \hline
	
	OpenTable &
	Нет &
	Да &
	Нет &
	Нет \\ \hline
	
	Poster POS &
	Да &
	Частично &
	Нет &
	Да \\ \hline
	
	Разрабатываемое приложение &
	Да &
	Да &
	Да &
	Да \\ \hline
\end{tabular}

\end{table}

Проведенный анализ показал, что существующие решения преимущественно специализируются на отдельных аспектах ресторанного бизнеса. Сервисы доставки ориентированы на оформление заказов и логистику, тогда как платформы бронирования позволяют только резервировать столики без интеграции с системой предзаказа блюд.

Кроме того, большинство решений не предоставляют пользователю возможности одновременно выбрать столик, оформить предварительный заказ и указать параметры посещения заведения. Отсутствие единой системы взаимодействия между клиентами, официантами и администраторами приводит к необходимости использования нескольких независимых сервисов.

Разрабатываемое веб-приложение ориентировано на устранение указанных ограничений и объединяет следующие возможности:

\begin{itemize}
\item онлайн-бронирование столиков с использованием интерактивной карты заведения;

\item предварительный заказ блюд к моменту посещения;

\item оформление заказов с доставкой и навынос;

\item управление заказами и бронированиями в единой системе;

\item инструменты администрирования персонала и столиков;

\item разграничение ролей пользователей (посетитель, официант, администратор).

\end{itemize}

Таким образом, разработка комплексного веб-приложения для сети кафе является актуальной задачей, направленной на повышение удобства взаимодействия клиентов с заведением и оптимизацию внутренних бизнес-процессов.

\subsection{Анализ программных инструментов разработки}

Для реализации проекта рассматриваются современные технологии веб-разработки, обеспечивающие высокую производительность, масштабируемость и удобство сопровождения приложения.

\subsubsection{Анализ технологий клиентской части}

В качестве технологий разработки пользовательского интерфейса были рассмотрены React + Next.js, Vue.js и Angular. Сравнение представлено в таблице~\ref{tab:frontend}.

\begin{table}[ht!]
\centering
\caption{Сравнение frontend-технологий}
\label{tab:frontend}
\begin{tabular}{|
>{\raggedright\arraybackslash}p{3cm}|
>{\raggedright\arraybackslash}p{5.2cm}|
>{\raggedright\arraybackslash}p{4.2cm}|}
\hline
\multicolumn{1}{|c|}{\textbf{Технология}} &
\multicolumn{1}{c|}{\textbf{Преимущества}} &
\multicolumn{1}{c|}{\textbf{Недостатки}} \\ \hline
	React + Next.js &
	Высокая производительность, развитая экосистема, поддержка SSR, удобная маршрутизация, масштабируемость &
	Требует более сложной архитектурной организации проекта \\ \hline
	
	Vue.js &
	Простота изучения и низкий порог входа &
	Меньшая экосистема для крупных проектов \\ \hline
	
	Angular &
	Полноценный framework со встроенными инструментами &
	Высокая сложность и избыточность для данного проекта \\ \hline
\end{tabular}

\end{table}

Для разработки клиентской части был выбран фреймворк Next.js, основанный на библиотеке React. Использование Next.js обеспечивает удобную организацию маршрутизации, поддержку серверного рендеринга, оптимизацию производительности и возможность построения масштабируемого пользовательского интерфейса.

Дополнительно используются следующие технологии:

\begin{itemize}
\item TypeScript --- для повышения типобезопасности приложения;

\item Tailwind CSS --- для ускорения разработки адаптивного пользовательского интерфейса;

\item Context API --- для управления состоянием авторизации и локализации;

\item App Router --- для организации структуры маршрутов приложения.

\end{itemize}

\subsubsection{Анализ технологий серверной части}

Для реализации серверной части были рассмотрены Express.js, NestJS и Spring Boot. Сравнение представлено в таблице~\ref{tab:backend}.

\begin{table}[ht!]
\centering
\caption{Сравнение backend-технологий}
\label{tab:backend}
\begin{tabular}{|
>{\raggedright\arraybackslash}p{3cm}|
>{\raggedright\arraybackslash}p{5.2cm}|
>{\raggedright\arraybackslash}p{4.2cm}|}
\hline
\multicolumn{1}{|c|}{\textbf{Технология}} &
\multicolumn{1}{c|}{\textbf{Преимущества}} &
\multicolumn{1}{c|}{\textbf{Недостатки}} \\ \hline

	NestJS &
	Модульная архитектура, поддержка TypeScript, встроенные механизмы DI и Guards &
	Более высокий порог входа \\ \hline
	
	Express.js &
	Минималистичность и простота &
	Требует самостоятельной организации архитектуры \\ \hline
	
	Spring Boot &
	Высокая надежность и производительность &
	Высокая сложность и ресурсоемкость \\ \hline
\end{tabular}

\end{table}

В результате анализа для реализации серверной части был выбран фреймворк NestJS, работающий на платформе Node.js. Данный выбор обусловлен следующими преимуществами:

\begin{itemize}
\item поддержка модульной архитектуры;

\item встроенные механизмы аутентификации и авторизации;

\item поддержка TypeScript;

\item удобная организация бизнес-логики;

\item возможность масштабирования проекта;

\item поддержка REST API.

\end{itemize}

\subsubsection{Выбор системы управления базами данных}

Для хранения данных рассматривались PostgreSQL и MongoDB.

MongoDB удобна при работе с неструктурированными данными, однако для системы бронирования столиков и управления заказами важна строгая структура данных, поддержка транзакций и наличие связей между сущностями.

В качестве системы управления базами данных была выбрана PostgreSQL, обеспечивающая:

\begin{itemize}
\item надежное хранение структурированных данных;

\item поддержку транзакций;

\item высокую производительность;

\item поддержку сложных SQL-запросов;

\item масштабируемость и стабильность.

\end{itemize}

\subsubsection{Выбор ORM}

Для взаимодействия с базой данных используется Prisma ORM.

Основными преимуществами Prisma являются:

\begin{itemize}
\item типобезопасность запросов;

\item автоматическая генерация клиента базы данных;

\item поддержка миграций;

\item снижение вероятности ошибок при работе с SQL-запросами;

\item удобная интеграция с TypeScript и NestJS.
\end{itemize}

Таким образом, выбранный стек технологий (Next.js, React, TypeScript, NestJS, PostgreSQL и Prisma ORM) обеспечивает высокую производительность, масштабируемость и удобство сопровождения веб-приложения для сети кафе.


\subsection{Формулировка цели и задач работы}

На основе проведенного анализа существующих решений (см. подраздел 1.1) и выбранных технологий (см. подраздел 1.2) была сформулирована цель разработки.

\textbf{Цель работы} --- разработка веб-приложения для управления заказами, бронированием столиков и предзаказом блюд в сети кафе.

Для достижения поставленной цели необходимо решить следующие \textbf{задачи}:

\begin{enumerate}
	\item Провести анализ существующих решений в области онлайн-заказа еды и автоматизации ресторанного бизнеса;
	
	\item Разработать архитектуру клиент-серверного приложения и спроектировать структуру базы данных для хранения информации о пользователях, заказах, блюдах и бронированиях;
	
	\item Реализовать систему аутентификации и авторизации пользователей с разграничением прав доступа (пользователь, администратор, официант);
	
	\item Разработать функционал оформления заказов, включая различные типы заказов (в заведении, доставка, на вынос);
	
	\item Реализовать систему бронирования столиков с возможностью выбора времени и краткосрочной брони;
	
	\item Разработать механизм предзаказа блюд с учетом времени посещения заведения;
	
	\item Реализовать систему обработки заказов персоналом и управление статусами заказов;
	
	\item Провести тестирование разработанного веб-приложения и оценить его работоспособность.
\end{enumerate}